\section{Теоретико-множественные операции над языками}
\label{sec:SetOperations}
\tikzsetfigurename{FLT_SetOps_}

Так как язык~--- это \emph{множество} слов, то над языками естественным образом определены теоретико-множественные операции, такие как объединение, пересечение, дополнение.
\begin{itemize}
    \item $L_1 \cup L_2 = \{ w \mid w \in L_1 \text{ или } w \in L_2\}$
    \item $L_1 \cap L_2 = \{ w \mid w \in L_1 \text{ и } w \in L_2\}$
    \item $\overline{L} = \{ w \mid w \in \Sigma^* \text{ и } w \notin L\} = \Sigma^* \setminus L$, где $L$~--- язык над алфавитом $\Sigma$ .
\end{itemize}

Многие прикладные задачи, в которых возникают формальные языки, достаточно естественным образом формулируются в теоретико-множественных терминах.
Так, задача распознавания~--- это задача проверки принадлежности элемента множеству.
Далее мы рассмотрим прикладные задачи, решение которых требует, например, проверки включения одного языка в другой, или проверки непустоты пересечения двух языков.
Разрешимость таких задач, алгоритмы решения, их сложность и другие свойства, зависят от свойств языков.
Это даёт дополнительную связь теоретико-языковых конструкций с решением прикладных задач, позволяя более аккуратно рассуждать о свойствах получаемых решений.
